\documentclass[10pt,a4paper]{amsart}
\usepackage[margin=19mm]{geometry}
\usepackage{amsmath,amssymb,mathtools}
\usepackage[scr=rsfs,cal=euler]{mathalfa}
\usepackage{xfrac}
\usepackage[unicode,hidelinks,bookmarks=false]{hyperref}
\newcommand{\ie}{i.e.\ }
\newcommand{\cf}{cf.\ }
\newcommand{\resp}{respectively\ }
\newcommand{\Subsection}{\subsection}
\providecommand{\nobreakdash}{\nobreak\hbox{-}}
\setlength{\emergencystretch}{2em}
\multlinegap0pt
\usepackage{mathtools}
\usepackage{booktabs}
\usepackage{array}
\usepackage{tikz}
\usetikzlibrary{arrows.meta,positioning,shapes.geometric}
\numberwithin{equation}{section}
\newtheorem{theorem}{Theorem}[section]
\newtheorem{lemma}[theorem]{Lemma}
\newtheorem{proposition}[theorem]{Proposition}
\newtheorem{conjecture}[theorem]{Conjecture}
\newtheorem{definition}[theorem]{Definition}
\newtheorem{question}[theorem]{Question}
\newtheorem*{proposition*}{Proposition}
\newtheorem{remark}[theorem]{Remark}
\newtheorem{example}[theorem]{Example}
\DeclareMathOperator{\Sym}{Sym}
\DeclareMathOperator{\Gr}{Gr}
\DeclareMathOperator{\rk}{rk}
\newcommand{\C}{\mathbb C}
\newcommand{\Q}{\mathbb Q}
\newcommand{\Z}{\mathbb Z}
\renewcommand{\P}{{\mathbb P}}
\DeclareMathOperator{\pol}{Pol}
\DeclareMathOperator{\Pol}{Pol}
\DeclareMathOperator{\Pl}{Pl}
\newcommand{\Dder}{\mathfrak D}
\newcommand{\NN}{\operatorname{NN}}
\newcommand{\stir}[2]{\genfrac{[}{]}{0pt}{}{#1}{#2}}
\newcommand{\StirTwo}[2]{\genfrac{\{} {\}}{0pt}{}{#1}{#2}}
\begin{document}
\title[Positivity in classical enumerative geometry: a case study i]{Positivity in classical enumerative geometry: a case study in synchronized AI-assisted mathematics}
\author{Gergely Bérczi; László M. Fehér}
\date{}
\maketitle
\noindent\emph{Source and licence.} Positivity in classical enumerative geometry: a case study in synchronized AI-assisted mathematics. Version arXiv:2605.25271v1 (CC BY 4.0). This material is distributed under the Creative Commons Attribution 4.0 International licence, http://creativecommons.org/licenses/by/4.0/, as recorded in the arXiv OAI licence metadata for this exact version and shown on the abstract page https://arxiv.org/abs/2605.25271v1. Full rights notice: licenses/arxiv-2605.25271-CC-BY-4.0.txt.\par\smallskip
\noindent\textit{Editorial scope.} Counted unit: the original \texttt{theorem} environment \texttt{thm:NN-binomiality} at source lines 638--656 together with its complete proof at lines 658--697; the delivered document keeps source lines 351--818 so that every referenced label and every citation resolves inside it. Native editable source is the paper's own concatenated TeX; the AMS wrapper is editorial formatting only. No publisher class, upstream script or unrelated artwork is redistributed.
\section{A binomiality theorem for symmetric powers}
This section describes the first mathematical phase of the workflow.  The starting point was the closed formula for $c_2(n,d)$ in Proposition~\ref{prop:c2}, first suggested experimentally by AlphaEvolve.  That formula exhibited exactly the binomial structure predicted by Feh\'er--Juh\'asz and led to the search for a general pattern for $c_3(n,d)$ and $c_4(n,d)$.  These patterns were then used as input for ChatGPT~5.5 Pro, which provided the factorial-moment
identity of Lemma~\ref{lem:factorial-moments} and the proof of the strong
binomiality Theorem \ref{thm:strong-binomiality}. The proof converts sums over weak compositions into binomial coefficients, and hence turns the Chern-root power sums into polynomials in the quantities $T_r(n,d)$. We then study the K-theory analogue where GPT Pro~5.5 successfully generalised Theorem \ref{thm:strong-binomiality} to the closed formula of Theorem \ref{thm:NN-binomiality}, whose diagonal coefficients recover the original Chern classes.

Let
\[
  A_{n,d}:=\{\alpha\in\Z_{\ge0}^n:|\alpha|=d\},\qquad
  y_\alpha:=\alpha_1x_1+\cdots+\alpha_nx_n.
\]
Then recall from the Introduction the Chern classes
\[
  \prod_{\alpha\in A_{n,d}}(1+y_\alpha)=\sum_{k\ge0}c_k(n,d).
\]
A two-parameter expression need not be polynomial in both parameters (also called separately polynomial function), even if it has a simple closed form.  For example, $d^n$ is a polynomial in $d$ for every fixed $n$, but after setting $d=2$ one obtains $2^n$, which is not a polynomial in $n$.  Thus polynomiality in the rank $n$ and the degree $d$ is already a genuine structural constraint.

\begin{conjecture}[Separate polynomiality, \cite{feher2025polynomiality}]\label{conj:polynomiality}
For each fixed $k$ and each partition $\lambda\vdash k$, the coefficient $f_\lambda(n,d)$ in
\[
  c_k(n,d)=\sum_{\lambda\vdash k}f_\lambda(n,d)e_\lambda
\]
is polynomial in $d$ for fixed $n$, and polynomial in $n$ for fixed $d$.
\end{conjecture}

The following conjecture is a stronger form of this polynomiality statement: it predicts not only polynomial dependence on the parameters separately, but polynomial dependence through a distinguished family of binomial coefficients.

\begin{conjecture}[Strong binomial form, \cite{feher2025polynomiality}]\label{conj:binomial-form}
For each fixed $k$ and each partition $\lambda\vdash k$, the coefficient $f_\lambda(n,d)$ in
\[
  c_k(n,d)=\sum_{\lambda\vdash k}f_\lambda(n,d)e_\lambda
\]
is a polynomial in finitely many binomial coefficients of the form $\binom{d+n+a}{n+b}$.
\end{conjecture}

The first nontrivial low-degree case already exhibits the special binomial structure.

\begin{proposition}[Closed formula for $c_2(n,d)$]\label{prop:c2}
One has
\[
  c_2(n,d)=\binom{d+n}{n+1}e_2+
  \left(
    \frac12\binom{d+n-1}{n}^2
    -\frac12\binom{d+n-1}{n}
    -\binom{d+n-1}{n+1}
  \right)e_1^2.
\]
This formula was  found experimentally by AlphaEvolve, recovering the result of \cite{ciliberto_zaidenberg2020onfanoschemes}.
\end{proposition}

\begin{proof}
Let
\[
  u=\binom{d+n-1}{n},\qquad v=\binom{d+n}{n+1}.
\]
For $m\ge1$ put
\[
  P_m(n,d):=\sum_{\alpha\in A_{n,d}}y_\alpha^m.
\]

A direct generating-function calculation gives
\[
  P_1(n,d)=\sum_{\alpha\in A_{n,d}}y_\alpha=u\,e_1
\]
and
\[
  P_2(n,d)=\sum_{\alpha\in A_{n,d}}y_\alpha^2=(2v-u)e_1^2-2v\,e_2.
\]
Newton's identity $2c_2=P_1^2-P_2$ now yields the formula.
\end{proof}


It turns out that the much stronger version of Conjecture \ref{conj:binomial-form} holds: the coefficients of the Chern classes are polynomials of an explicit short list of binomial coefficients $T_r(n,d)$ below. 

\begin{theorem}[Strong binomiality]\label{thm:strong-binomiality}
For every $k\ge0$, there is a universal polynomial
\[
  Q_k\in\Q[u_1,\ldots,u_k,z_1,\ldots,z_k]
\]
such that, for all $n,d$,
\[
  c_k(n,d)=Q_k\bigl(T_1(n,d),\ldots,T_k(n,d);e_1,\ldots,e_k\bigr),
\]
where
\[
  T_r(n,d):=\binom{d+n-1}{n+r-1},\qquad 1\le r\le k.
\]
Equivalently, $f_\lambda(n,d)\in\Q[T_1(n,d),\ldots,T_k(n,d)]$ for every $\lambda\vdash k$.  In particular, Conjecture \ref{conj:binomial-form} holds.
\end{theorem}

\begin{lemma}[Factorial moments of weak compositions]\label{lem:factorial-moments}
Let $i_1,\ldots,i_s$ be distinct indices and let $m_1,\ldots,m_s$ be positive integers.  Put $M=m_1+\cdots+m_s$ and write $a^{\underline m}=a(a-1)\cdots(a-m+1)$.  Then
\begin{equation}\label{eq:factorial-moments}
  \sum_{\alpha\in A_{n,d}}\prod_{j=1}^s \alpha_{i_j}^{\underline {m_j}}
  =
  \left(\prod_{j=1}^s m_j!\right)\binom{d+n-1}{n+M-1}.
\end{equation}
\end{lemma}

\begin{proof}
Taking generating functions in $d$ gives
\[
  \sum_{d\ge0}\left(\sum_{\alpha\in A_{n,d}}\prod_{j=1}^s \alpha_{i_j}^{\underline {m_j}}\right)t^d
  =
  \left(\prod_{j=1}^s\frac{m_j!t^{m_j}}{(1-t)^{m_j+1}}\right)(1-t)^{-(n-s)}.
\]
This is
\[
  \left(\prod_{j=1}^s m_j!\right)t^M(1-t)^{-(n+M)}.
\]
Taking the coefficient of $t^d$ proves the lemma.
\end{proof}

\begin{proof}[Proof of Theorem \ref{thm:strong-binomiality}]
Fix $m\ge1$.  We first show that
\begin{equation}\label{eq:P-m-in-ring}
  P_m(n,d)\in \Z[T_1(n,d),\ldots,T_m(n,d)]\otimes_{\Z}\Lambda_m,
\end{equation}
where $\Lambda_m$ denotes homogeneous symmetric functions of degree $m$.

Let $\mu=(\mu_1,\ldots,\mu_s)\vdash m$.  The coefficient of a monomial of type $\mu$ in $P_m$ is a multinomial factor times
\[
  M_\mu(n,d):=\sum_{\alpha\in A_{n,d}}\alpha_1^{\mu_1}\cdots\alpha_s^{\mu_s}.
\]
Using Stirling numbers of the second kind,
\[
  a^q=\sum_{r=1}^{q}\StirTwo{q}{r}a^{\underline r},
\]
and applying Lemma \ref{lem:factorial-moments}, we obtain
\begin{equation}\label{eq:mixed-moment}
  M_\mu(n,d)=
  \sum_{1\le r_j\le \mu_j}
  \left(\prod_{j=1}^s \StirTwo{\mu_j}{r_j}r_j!\right)
  T_{r_1+\cdots+r_s}(n,d).
\end{equation}
Thus each monomial-symmetric coefficient of $P_m$ lies in $\Z[T_1,\ldots,T_m]$.  Since the elementary products form a basis in each degree over $\Q$, $P_m$ can be written universally as a polynomial in $T_1,\ldots,T_m$ and $e_1,\ldots,e_m$.

Newton identities for the roots $y_\alpha$ give
\[
  k c_k=\sum_{j=1}^{k}(-1)^{j-1}c_{k-j}P_j,\qquad c_0=1.
\]
By induction on $k$, the right hand side lies in $\Q[T_1,\ldots,T_k][e_1,\ldots,e_k]$.  Dividing by $k$ proves the theorem.
\end{proof}
\begin{example} \label{ex:small-k-cknd} For small $k$ we have
\[\begin{aligned}
c_1(n,d)&=T_1e_1,\\[3pt]
c_2(n,d)&=
\frac{T_1^2-T_1-2T_2}{2}e_{1}^2
+(T_1+T_2)e_2,\\[3pt]
c_3(n,d)&=
\frac{T_1^3-3T_1^2-6T_1T_2+2T_1+12T_2+12T_3}{6}e_{1}^3\\
&\quad+
\left(T_1^2+T_1T_2-T_1-5T_2-4T_3\right)e_1e_2\\
&\quad+
\left(T_1+3T_2+2T_3\right)e_3.
\end{aligned}\]
\end{example}
\medskip
\begin{remark} \label{strong-binom4schur}
Theorem\ref{thm:strong-binomiality} immediately implies that the Schur coefficents of $c_k(n,d)$ are also polynomials in the binomial expressions $T_r$.
\end{remark}

\subsection{The case of $d=2$: the Laksov--Lascoux--Thorup formula}
\smallskip

For $d=2$,
\[
  C_n(x):=c(\Sym^2\C^n)=\prod_{i=1}^n(1+2x_i)\prod_{1\le i<j\le n}(1+x_i+x_j).
\]

The case of $d=2$ can be calculated from \eqref{eq:mixed-moment} but ChatGPT~5.5 Pro gives the following shorter proof:



\begin{theorem}[ \cite{LLT} The $d=2$ determinant]\label{thm:d2det}
Let $\lambda\vdash k$ with $\ell(\lambda)\le n$.  Then
\[
  c_k(n,2)=\sum_{\lambda\vdash k,\ \ell(\lambda)\le n}A_\lambda(n)s_\lambda(x),
\]
where
\begin{equation}\label{eq:d2-coeff}
  A_\lambda(n)=2^{k-\binom n2}
  \det\left[\binom{2n-2j+1}{\lambda_i+n-i}\right]_{1\le i,j\le n}.
\end{equation}
\end{theorem}

\begin{proof}
Let $h_r$ be the complete symmetric functions, with $h_r=0$ for $r<0$.  For a strictly increasing sequence $I=(i_1<\cdots<i_n)$ put
\[
  s_I=\det(h_{i_q-p+1})_{1\le p,q\le n}.
\]
Set $A_{r,q}=2^r\binom{2n-2q+1}{r}$ and $B_{p,r}=h_{r-p+1}$.  By Cauchy--Binet,
\[
  \det(BA)=\sum_I s_I\det(A_I).
\]
Using
\[
  h_m=\sum_{a=1}^n \frac{x_a^{m+n-1}}{\prod_{b\ne a}(x_a-x_b)},
\]
one obtains
\[
  (BA)_{p,q}=\sum_{a=1}^n
  \frac{x_a^{n-p}(1+2x_a)^{2n-2q+1}}{\prod_{b\ne a}(x_a-x_b)}.
\]
Taking determinants,
\[
  \det(BA)=(-1)^{\binom n2}
  \frac{\det((1+2x_a)^{2n-2q+1})_{a,q}}{\prod_{a<b}(x_a-x_b)}.
\]
With $u_a=1+2x_a$,
\[
  \det(u_a^{2n-2q+1})_{a,q}
  =\left(\prod_a u_a\right)\prod_{a<b}(u_a^2-u_b^2).
\]
Since $u_a^2-u_b^2=4(x_a-x_b)(1+x_a+x_b)$, we get
\[
  \det(BA)=(-1)^{\binom n2}2^{n(n-1)}C_n(x).
\]
Reindexing $I$ by partitions gives \eqref{eq:d2-coeff}.
\end{proof}

\begin{example}
For small $k$,
\[
  c_1(n,2)=(n+1)s_1,
\]
\[
  c_2(n,2)=\frac{(n-1)(n+2)}2s_2+\frac{(n+1)(n+2)}2s_{11},
\]
\[
  c_3(n,2)=\frac{(n-2)(n-1)(n+3)}6s_3+
  \frac{(n+2)(n^2+n-3)}3s_{21}
  +\frac{(n+1)(n+2)(n+3)}6s_{111}.
\]
\end{example}

\medskip
\subsection{A $K$-theoretic generalization.}
The following normalized multiplicative generalization
was proved in this project by GPT Pro~5.5.  It is useful to regard it as a $K$-theoretic or finite-difference refinement of Theorem~\ref{thm:strong-binomiality}: the variables $X_i$ are multiplicative Chern-root variables, while the substitution $X_i=1-\zeta u_i$ extracts the additive Chern-root information order by order in $\zeta$.

Let
\[
  M_{n,d}(y,X)=\prod_{\alpha\in A_{n,d}}(1+yX^\alpha),
  \qquad X^\alpha=X_1^{\alpha_1}\cdots X_n^{\alpha_n}.
\]

This is the $K$-theoretic total Chern class, or equivariant motivic Chern class of the representation $\Sym^d(\mathbb C^n)$. Using a substitution we can compare it with the cohomological Chern class:

Define
\[
  \NN_{n,d}(z,\zeta,u):=
  (1+z)^{|A_{n,d}|}
  M_{n,d}\left(-\frac{z}{1+z},\,X_j\mapsto 1-\zeta u_j\right).
\]
Equivalently,
\begin{equation}\label{eq:NN-def}
  \NN_{n,d}(z,\zeta,u)=
  \prod_{\alpha\in A_{n,d}}
  \left(1+z\left(1-\prod_{j=1}^n(1-\zeta u_j)^{\alpha_j}\right)\right).
\end{equation}
Then the cohomological Chern class is the diagonal:
\begin{proposition}[Diagonal recovery of the Chern classes]\label{prop:NN-diagonal}
For every $k\ge0$,
\[
  [z^k\zeta^k]\NN_{n,d}(z,\zeta,u)=c_k(n,d).
\]
\end{proposition}

\begin{proof}
For each $\alpha\in A_{n,d}$,
\[
  1-\prod_{j=1}^n(1-\zeta u_j)^{\alpha_j}
  =\zeta\sum_{j=1}^n\alpha_j u_j+O(\zeta^2).
\]
To obtain $z^k\zeta^k$, one must choose the $z$-term from exactly $k$ factors in \eqref{eq:NN-def}.  Since each chosen factor contributes at least one power of $\zeta$, the total $\zeta$-degree $k$ forces us to take only the linear part displayed above.  Therefore
\[
  [z^k\zeta^k]\NN_{n,d}
  =\sum_{\substack{S\subseteq A_{n,d}\\ |S|=k}}
  \prod_{\alpha\in S}\left(\sum_j\alpha_j u_j\right),
\]
which is precisely the degree-$k$ homogeneous component of
\[
  \prod_{\alpha\in A_{n,d}}
  \left(1+\sum_j\alpha_j u_j\right).
\]
\end{proof}
 Using this we can see that the following theorem generalizes Theorem \ref{thm:strong-binomiality}:

\begin{theorem}[Binomiality for the $\NN$-analogue]\label{thm:NN-binomiality}
For every fixed pair $q,m\ge0$, the coefficient
\[
  [z^q\zeta^m]\NN_{n,d}(z,\zeta,u)
\]
is a symmetric homogeneous polynomial of degree $m$ in $u_1,\ldots,u_n$.  If
\[
  [z^q\zeta^m]\NN_{n,d}
  =\sum_{\lambda\vdash m}F_{q,m,\lambda}(n,d)e_\lambda(u),
\]
then
\[
  F_{q,m,\lambda}(n,d)\in\Q[T_1(n,d),\ldots,T_m(n,d)].
\]
Moreover,
\[
  [z^q\zeta^m]\NN_{n,d}=0\qquad\text{if }q>m.
\]
\end{theorem}

\begin{proof}
For $\alpha\in A_{n,d}$ put
\[
  A_\alpha(\zeta,u):=1-\prod_{j=1}^n(1-\zeta u_j)^{\alpha_j}.
\]
Then \eqref{eq:NN-def} says
\[
  \NN_{n,d}=\prod_{\alpha\in A_{n,d}}(1+zA_\alpha),
\]
and hence
\[
  [z^q]\NN_{n,d}=e_q(A_\alpha:\alpha\in A_{n,d}).
\]
Newton identities express this elementary symmetric function universally in the power sums
\[
  P_s^{\NN}:=\sum_{\alpha\in A_{n,d}}A_\alpha^s,
  \qquad 1\le s\le q.
\]
Now
\[
  A_\alpha(\zeta,u)=
  \sum_{\beta\ne0}(-1)^{|\beta|+1}
  \binom{\alpha}{\beta}\zeta^{|\beta|}u^\beta,
  \qquad
  \binom{\alpha}{\beta}:=\prod_{j=1}^n\binom{\alpha_j}{\beta_j}.
\]
Thus the coefficient of $\zeta^m u^\mu$ in $P_s^{\NN}$ is a finite linear combination of sums of the form
\[
  \sum_{\alpha\in A_{n,d}}
  \prod_j
  \binom{\alpha_j}{\beta_j^{(1)}}
  \binom{\alpha_j}{\beta_j^{(2)}}\cdots
  \binom{\alpha_j}{\beta_j^{(s)}}.
\]
For fixed $j$, the product of binomial coefficients is a polynomial in $\alpha_j$ and can be expanded in the falling-factorial basis.  The total falling-factorial degree is at most $m$.  Applying Lemma~\ref{lem:factorial-moments} shows that every such sum lies in
\[
  \Q[T_1(n,d),\ldots,T_m(n,d)].
\]
Therefore the same is true for every coefficient of each $P_s^{\NN}$, and Newton identities imply the assertion for $[z^q]\NN_{n,d}$.  Extracting $\zeta^m$ proves the coefficient statement.  Finally, each $A_\alpha$ has $\zeta$-order at least one, so a term of $z$-degree $q$ has $\zeta$-order at least $q$.  Hence $[z^q\zeta^m]\NN_{n,d}=0$ for $q>m$.
\end{proof}




\begin{example}[Small $\NN$-coefficients]\label{ex:NN-small}
Write $T_r=T_r(n,d)$.  The first coefficients are
\[
  [z\zeta]\NN_{n,d}=T_1e_1,
\]
\[
  [z\zeta^2]\NN_{n,d}=-T_2e_{1,1}+T_2e_2,
\]
\[
  [z^2\zeta^2]\NN_{n,d}=
  \frac{T_1^2-T_1-2T_2}{2}e_{1,1}+(T_1+T_2)e_2,
\]
\[
  [z\zeta^3]\NN_{n,d}=T_3e_{1,1,1}-2T_3e_{2,1}+T_3e_3,
\]
\begin{align*}
  [z^2\zeta^3]\NN_{n,d}={}&
  (3T_3+2T_2-T_1T_2)e_{1,1,1}\\
  &+(T_1T_2-6T_3-5T_2)e_{2,1}
  +(3T_3+3T_2)e_3,
\end{align*}
\begin{align*}
  [z^3\zeta^3]\NN_{n,d}={}&
  \frac{T_1^3-3T_1^2-6T_1T_2+12T_3+12T_2+2T_1}{6}e_{1,1,1}\\
  &+(T_1^2+T_1T_2-4T_3-5T_2-T_1)e_{2,1}
  +(2T_3+3T_2+T_1)e_3.
\end{align*}

Notice that the coefficient of $[z^k\zeta^k]$ agrees with $c_k(n,d)$ in \ref{ex:small-k-cknd}, as the diagonal recovery Proposition \ref{prop:NN-diagonal} claims.
\end{example}

\subsection{Leading terms and applications}
Formula \eqref{eq:mixed-moment} immediately implies that the $d$ leading coefficients of $c_k(n,d)$ in the elementary symmetric basis are
\[ [e_\lambda]{c_k(n,d)}
  =
  \frac{1}{\prod_{i\ge1}m_i!\,((d-1)!)^{\ell(\lambda)}}
  n^{(d-1)\ell(\lambda)}
  +O\left(n^{(d-1)\ell(\lambda)-1}\right).
\]

This was conjectured in \cite[Conj 3.10]{feher2025polynomiality}. 

The leading terms for the Schur basis were obtained \cite{feher2025polynomiality}. That result led to the generalization of a theorem of Manivel of the degree of hypersurfaces containing a projective  $r$-space. Interestingly, the leading terms for the Schur basis are easier to calculate and their degree is uniform $n|\lambda|$. For the elementary symmetric basis the degrees can be much smaller due to a surprising cancellation.

The result on the leading coefficients of $c_k(n,d)$ in the elementary symmetric basis is more subtle, but we are not aware of any enumerative applications.

\subsection{Euler classes}
The Schur coefficients of the top Chern class or Euler class 

\begin{equation}\label{eq:etop}
  E_{n,d}(x):=c_{\rm top}(n,d)=
  \prod_{\substack{\alpha\in\Z_{\ge0}^n\\|\alpha|=d}}
  (\alpha_1x_1+\cdots+\alpha_nx_n).
\end{equation}
also has enumerative applications as it is explained e.g. in \cite{feher2025polynomiality}.  For $n=2$ let 
\[
c_{\mathrm{top}}(2,d)
=
\sum_{b=0}^{\left\lfloor (d+1)/2\right\rfloor}
\gamma_{d,b}\,
s_{(d+1-b,b)}(x_1,x_2).
\]
be the Schur expansion of the Euler class. Then ChatGPT 5.5 Pro gives

\[
\gamma_{d,b}
=
\sum_{r=b-1}^{d}
(-1)^{r-b}
\binom{r+1}{b}
d^{\,d+1-r}
\left[{d+1\atop d+1-r}\right],
\]

which is somewhat simpler than the formula of \cite{feher2025polynomiality}. This simplifies the formula of \cite{feher2025polynomiality} for the Euler characteristics of Fano varieties of lines of hypersurfaces. Here square brackets denote the unsigned Stirling numbers of the first kind.
For higher $n$ ChatGPT 5.5 Pro did not provide closed formulas.
\bigskip
\begin{remark}
Similar enumerative applications can be given for the $K$-theory Chern class, leading to calculation of Hilbert polynomials and Todd classes of Fano varieties. These applications will be published elsewhere.
\end{remark}

\section{A rank-two binomial refinement}

This section introduces the refined rank-two phenomenon that became the central target of the later AI-assisted proof stage.
We study the $n=2$, rank two case and write $c_k(2,d)$ in the Schur basis:
\[
  c_k(2,d)=\sum_{0\le j\le\lfloor k/2\rfloor}A_{k,j}(d)s_{(k-j,j)}(x_1,x_2).
\]
We then change coordinates once more by expanding each Schur coefficient $A_{k,j}(d)$ in the binomial basis of the degree parameter $d$:
\[
  A_{k,j}(d)=\sum_{r\ge0} B_{k,j,r}\binom{d}{r}.
\]
This change of basis, which was previously unknown in the literature and was motivated by Theorem \ref{thm:strong-binomiality}, revealed positivity and log-concavity patterns that were invisible in the original $d$-power basis.  

\begin{definition}\label{def:binom} We will call a polynomial
\[
  f(d)=\sum_{r\ge0} b_r\binom{d}{r}
\]
\emph{binomially positive} if $b_r\ge0$ for all $r$.  We will call it \emph{binomially log-concave} if the nonzero sequence $(b_r)_r$ has no internal zeroes and is log-concave, that is, $b_r^2\ge b_{r-1}b_{r+1}$.
\end{definition}
The following rank-two conjectures predict that, for each fixed $k$ and $j$, the Schur coefficient $A_{k,j}(d)$ is binomially positive and binomially log-concave as a polynomial in $d$.  


\begin{conjecture}[Rank-two binomial positivity]\label{conj:A}
For all $k,j,r$ with $0\le j\le\lfloor k/2\rfloor$, one has
\[
  B_{k,j,r}\ge0.
\]
Equivalently, for each fixed $k$ and $j$, the polynomial $A_{k,j}(d)$ is binomially positive.
\end{conjecture}

\begin{conjecture}[Rank-two binomial log-concavity]\label{conj:B}
For fixed $k$ and $j$, the nonzero sequence $(B_{k,j,r})_r$ is log-concave:
\[
  B_{k,j,r}^2\ge B_{k,j,r-1}B_{k,j,r+1}.
\]
Equivalently, $A_{k,j}(d)$ is binomially log-concave as a polynomial in $d$.
\end{conjecture}

\begin{thebibliography}{99}
\bibitem[LLT89]{LLT} Dan Laksov, Alain Lascoux, and Anders Thorup, \emph{On {Giambelli}'s theorem on complete correlations}, Acta Mathematica \textbf{162} (1989), no.~3--4, 143--199.
\bibitem[CZ20]{ciliberto_zaidenberg2020onfanoschemes} Ciro Ciliberto and Mikhail Zaidenberg, \emph{On {Fano} schemes of complete intersections}, Polynomial Rings and Affine Algebraic Geometry (Shigeru Kuroda, Nobuharu Onoda, and Gene Freudenburg, eds.), Springer Proceedings in Mathematics \& Statistics, vol. 319, Springer, Cham, 2020, pp.~1--39.
\bibitem[FJ25]{feher2025polynomiality} L{\'a}szl{\'o}~M. Feh{\'e}r and Andr{\'a}s~P. Juh{\'a}sz, \emph{Polynomiality of the {Stirling} coefficients of {$c(\mathrm{Pol}^d(\mathbb{C}^n))$} and {Fano} schemes}, arXiv preprint arXiv:2509.01725, 2025, Available at \url{https://arxiv.org/abs/2509.01725}.
\end{thebibliography}
\end{document}
